\chapter{Rustlets}

This chapter describes Rustlets from the application developer point of
view. It focuses on the functional execution model, the framework
surface presented to Rustlet authors, and the abstract isolation
principles enforced by the secure element.

\section{General Principle}

A Rustlet is an application payload compiled as an independent FAE image
and later executed inside the secure element under kernel mediation.
From the developer point of view, a Rustlet is not a hosted process and
not a kernel extension. It is an application component that exposes a
small binary interface to the kernel and that is invoked through APDU
selection and dispatch.

The current model is stateful. A Rustlet is associated with an
application state that is created during installation and then reused
across later APDU processing calls. Selection activates a Rustlet
instance and initializes its working data for execution. In other words,
selection establishes which Rustlet instance becomes the current
execution context, while installation is responsible for creating the
resident Rustlet state itself.

This distinction matters for the long-term architecture:
\begin{itemize}
\item \texttt{SELECT} determines which Rustlet instance is active on the
  current command path;
\item \texttt{INSTALL} creates or configures one application instance;
\item \texttt{process\_apdu()} consumes and updates the active
  in-memory state of that instance.
\end{itemize}

The present implementation currently keeps one selected Rustlet context
active at a time, but the conceptual model is broader. A given Rustlet
code image may give rise to multiple installed instances, each with its
own serialized state. The secure element therefore has to distinguish:
\begin{itemize}
\item the executable Rustlet code image;
\item the installed Rustlet instance;
\item the current selected instance used for command dispatch.
\end{itemize}

This is also why Rustlet state should be understood as a persistent
application object rather than as a transient stack frame. During
execution, a Rustlet may use temporary working data, heap allocations
and APDU-local buffers. Only a subset of this data belongs to the
resident Rustlet state and should survive across commands.

% The current prototype still keeps the resident state in RAM and does
% not yet serialize it after each APDU. The intended model is that the
% persistent subset of the Rustlet state is serialized at the end of each
% successful APDU processing step, so multiple installed instances can be
% stored and later reselected.

The effective realization of selection, installation, instance storage,
and activation is described from an implementation point of view in
Chapter~8 (Reference Implementation), in
Section~\ref{sec:impl-rustlet-runtime}.

\section{Framework and API}

\subsection{Framework}

The Rustlet framework is intentionally small. A Rustlet developer
defines:
\begin{itemize}
\item one Rustlet state type;
\item one installation function that creates this state;
\item one implementation of the \texttt{Rustlet} trait for that state;
\item one macro invocation exporting the Rustlet entry point expected by
  the runtime.
\end{itemize}

The state type represents the resident application state that survives
across APDU calls. The installation function is the constructor of that
resident state. The \texttt{Rustlet} trait supplies the command
processing behavior once an instance exists.

The declaration macro is:
\begin{verbatim}
declare_rustlet!(MyRustlet);
declare_rustlet!(MyRustlet, 2048usize);
declare_rustlet!(MyRustlet, 2048usize, install_custom);
\end{verbatim}

This macro exports the FAE \texttt{start()} entry point and binds the
Rustlet type, its heap reservation and its installation function to the
runtime. It therefore acts as the bridge between the high-level Rustlet
code written by the developer and the low-level ABI expected by the
kernel.

The \texttt{Rustlet} trait itself should be understood abstractly as the
interface of an installed Rustlet instance. Its main role is to process
one APDU within the current selected context and to update the resident
state accordingly. In the present implementation, its main entry point has
the following signature:
\begin{center}
\small\texttt{process\_apdu(\&mut self, \&mut RustletCtx) -> ApduStatus}
\end{center}
The same trait also exposes the runtime save/load boundary used for the
persistent subset of the resident Rustlet state.

The implementation details of the runtime adapters, macro expansion, FAE
entry point generation and selected-context descriptor are described in
Chapter~8 (Reference Implementation), in
Section~\ref{sec:impl-rustlet-runtime}.

\subsection{API}

The current Rustlet API is intentionally narrow and typed. The central runtime
object is \texttt{RustletCtx}, while ordinary Rustlet code interacts with
APDUs through the typed \texttt{Apdu<State>} interface exposed by the
runtime.

\texttt{RustletCtx} is the shared ABI region between the kernel and the
currently executing Rustlet. It carries:
\begin{itemize}
\item the ABI version;
\item the APDU header fields;
\item the status word;
\item one fixed-size APDU payload buffer;
\item one fixed-size state-serialization slice;
\item a small amount of direction and length metadata.
\end{itemize}

The current binary layout is fixed to one 512-byte shared page. It is split
into two 256-byte halves:

\begin{center}
\begin{tabular}{|l|l|p{7.0cm}|}
\hline
\textbf{Offset} & \textbf{Size} & \textbf{Role} \\
\hline
\texttt{0x000} & \texttt{256} &
Control and secondary buffer. On Rustlet entry this contains ABI metadata,
APDU header/status fields, flags, the APDU data length and serialized state. \\
\hline
\texttt{0x100} & \texttt{256} &
APDU payload buffer. Before outgoing mode it contains incoming command data;
after outgoing mode it contains response data. \\
\hline
\end{tabular}
\end{center}

The control and secondary half currently contains:
\begin{itemize}
\item \texttt{abi\_version: u32};
\item five command/status bytes, interpreted as
  \texttt{CLA INS P1 P2 P3} on entry and \texttt{SW1 SW2 00 00 00} on return;
\item one flags byte;
\item one APDU data-length byte;
\item one serialized-state length byte;
\item \texttt{244} bytes of serialized-state storage.
\end{itemize}

Three invariants follow from this layout:
\begin{itemize}
\item the shared page size remains exactly \texttt{512} bytes, so the kernel
  can map the Rustlet ABI as one small gate region;
\item the APDU payload buffer reserves \texttt{256} bytes, but the current
  short-APDU ABI accepts at most \texttt{255} useful payload bytes because the
  data length is represented by one byte;
\item the control and secondary half may be reused by kernel protocol code as
  temporary scratch space before entering a Rustlet, but it must be reset and
  reinitialized before Rustlet code observes the ABI page.
\end{itemize}

Its low-level functional meaning is simple:
\begin{itemize}
\item before outgoing mode is enabled, the payload buffer contains the
  incoming APDU data;
\item after outgoing mode is enabled, the same payload buffer becomes
  the outgoing response area;
\item at the end of processing, the Rustlet returns an
  \texttt{ApduStatus}, while the runtime owns the corresponding
  \texttt{SW1/SW2} staging.
\end{itemize}

However, this is \emph{not} the preferred programming model for normal
Rustlet code. The preferred model is the typed APDU state machine:
\begin{itemize}
\item \texttt{Apdu<Command>}, the initial command-processing state;
\item \texttt{Apdu<Receiving>}, after the incoming phase has been
  accepted;
\item \texttt{Apdu<Sending>}, after the outgoing phase has been
  declared.
\end{itemize}

This state machine provides a more precise expression of the APDU case
currently being processed. In practice:
\begin{itemize}
\item a Rustlet builds \texttt{Apdu::new(ctx)} at the beginning of
  \texttt{process\_apdu()};
\item it inspects the command header in the \texttt{Command} state;
\item it moves to \texttt{Receiving} only if the command needs incoming
  payload;
\item it moves to \texttt{Sending} only if the command produces outgoing
  payload.
\end{itemize}

This typed API is the Rustlet-side normalization point of the current
framework. It is intentionally distinct from:
\begin{itemize}
\item the kernel transport object, which remains responsible for T=0
  procedure bytes, \texttt{61xx}, \texttt{6Cxx}, and \texttt{GET
  RESPONSE};
\item the raw shared memory layout, which remains an ABI detail rather
  than the preferred application programming model.
\end{itemize}

The Rustlet-facing API also includes:
\begin{itemize}
\item \texttt{ApduHeader}, which exposes the decoded command header;
\item \texttt{ApduStatus}, which carries the result status returned to
  the kernel;
\item \texttt{RustletCtx}, which represents the current runtime context
  made visible to the Rustlet;
\item \texttt{RustletHeapRegion}, which describes the heap reservation
  associated with the Rustlet instance;
\item \texttt{SelectedAppDescriptor} and \texttt{SelectedAppVtable},
  which belong to the kernel/Rustlet ABI boundary rather than to
  everyday Rustlet business logic.
\end{itemize}

From the developer point of view, these objects define a deliberately
small application framework:
\begin{itemize}
\item one shared runtime context object;
\item one typed APDU state machine;
\item one resident state object;
\item one trait-based processing interface;
\item one kernel-mediated crypto interface.
\end{itemize}

The effective binary layout of these objects, as well as the runtime and
kernel-side marshalling rules, are described in Chapter~8 (Reference
Implementation), in Section~\ref{sec:impl-rustlet-runtime}. In other
words, the present chapter states the programming model, while
Chapter~8 gives the underlying ABI and runtime mechanics that make this
model concrete.

\section{Security Domains Dedicated Framework and API}
\label{sec:rustlets-security-domain-framework}

A Rustlet may also be declared as a user-land Security Domain. In that
case, it remains a normal selectable Rustlet from the APDU point of
view, but it also exposes a dedicated Security Domain interface to the
kernel. This is the mechanism intended to let a Rustlet implement
GlobalPlatform administrative authority without moving all management
logic into the kernel.

In Oxide SE, Security Domains are modeled as registry objects like
applications and load-file related objects. A Security Domain instance is
therefore represented by an instance AID together with administrative metadata,
rather than by a hidden side structure unrelated to the registry.

The declaration macro is:
\begin{verbatim}
declare_security_domain!(MySecurityDomain);
declare_security_domain!(MySecurityDomain, 2048usize);
declare_security_domain!(MySecurityDomain, 2048usize, install_custom);
\end{verbatim}

The declared type must implement both:
\begin{itemize}
\item \texttt{Rustlet}, for ordinary APDU dispatch through
  \texttt{process\_apdu()};
\item \texttt{RustletSecurityDomain}, for Security Domain management
  and secure-channel hooks invoked by the kernel proxy.
\end{itemize}

This dual role is deliberate. A Security Domain is, in the present
architecture, first a Rustlet and then an administrative authority. It is
addressable and selectable like an application, but it also carries
management and secure-channel responsibilities. The APDU-facing side is
used for commands explicitly routed to the Security Domain itself. The
dedicated interface is used when the kernel needs a policy decision or
secure-channel operation while processing platform management commands.

Registry-wise, a Security Domain instance behaves like any other managed
object: it has an instance AID, it may have a \texttt{parent\_sd\_aid}, and it
carries a serialized administrative state. The bootstrap administrative
authority is defined as the first Security Domain instance recorded without a
\texttt{parent\_sd\_aid}. This instance is seeded by the kernel during
bootstrap together with its initial privilege state.

The current Rustlet-side Security Domain interface mirrors the kernel
view of Security Domain authority. It is useful to separate it into
several functional blocks, because this is also how GlobalPlatform
administrative behavior is decomposed in practice:
\begin{itemize}
\item \textbf{application-like behavior}, still implemented through
  \texttt{Rustlet::process\_apdu()} when the Security Domain itself is
  selected;
\item \textbf{management hooks}, such as
  \texttt{install\_for\_load()},
  \texttt{install\_for\_install()}, \texttt{delete\_aid()},
  \texttt{get\_data()}, \texttt{may\_manage\_applet()}, and
  \texttt{may\_make\_selectable()};
\item \textbf{secure-channel hooks}:
  \begin{itemize}
  \item initialization through \texttt{initialize\_update()};
  \item authentication through \texttt{external\_authenticate()};
  \item session inspection through \texttt{current\_security\_level()} and
    \texttt{secure\_channel\_open()};
  \item MAC-length queries through \texttt{current\_mac\_len()};
  \item teardown through \texttt{reset\_secure\_channel()}.
  \end{itemize}
\item \textbf{secure-messaging hooks}, namely
  \texttt{unwrap\_command()} and \texttt{wrap\_response()}, which
  transform protected APDUs once the secure channel is open.
\item \textbf{privilege and key-management concerns}, which already have
  a first place in the interface (\texttt{privileges()}) and already
  cover one first \texttt{PUT KEY} subset, but still require richer GP
  work such as broader keyset policy, rotation semantics, and
  diversification.
\end{itemize}

At the current stage, the runtime already provides one common privilege
representation for Rustlet Security Domains. It decodes the complete
first GlobalPlatform privilege byte carried by
\texttt{INSTALL [for install]} and therefore distinguishes, at least,
the following privilege bits:
\begin{itemize}
\item Security Domain;
\item DAP Verification;
\item Delegated Management;
\item Card Lock;
\item Card Terminate;
\item Card Reset;
\item CVM Management;
\item Mandated DAP Verification.
\end{itemize}

The Rustlet-side framework groups management hooks around the decoded
privilege state. Its default policy covers loading, installation, deletion,
applet management, and transitions to the selectable state through the
corresponding hooks described above. These decisions can already derive a
first policy view from the decoded privilege state. This privilege state is part of the
standard persistent administrative state carried by a Rustlet-backed
Security Domain instance, together with a minimal lifecycle field.
Under the current implementation, \texttt{INSTALL [for install and make
selectable]} initializes that lifecycle to a selectable state, and
\texttt{GET DATA 9F70} is served from that runtime-managed administrative
state. The second and third privilege bytes remain preserved as encoded
state for later interpretation. Under the current policy, installation
already enforces a first non-escalation rule: a Security Domain instance
cannot grant privilege bits it does not itself hold. This first
non-escalation check is enforced by the kernel proxy before the Rustlet
Security Domain hook is called. The Rustlet may therefore refine a
management decision by rejecting it, but it cannot extend base
privileges that the kernel proxy has already refused. During bootstrap,
the kernel creates one implicit root Security Domain instance and grants
that bootstrap instance the complete administrative privilege set currently
modeled by the platform.

The same rule also structures delegation. A Security Domain instance may only
create descendant objects under its own authority, and a child Security Domain
instance may not receive more privileges than its parent instance already
holds. In this model, the delegation path is expressed through
\texttt{parent\_sd\_aid} in the registry rather than through ad hoc side
tables. For already existing objects, the same hierarchy constrains later
management: selection, deletion, and authority checks only succeed when the
target object belongs to the active Security Domain instance sub-tree.
Operationally, the registry itself is treated as the visibility boundary for
that administrative context: ordinary lookups are resolved from the point of
view of the Security Domain instance currently selected for the active secure
channel, so that objects outside the visible sub-tree are treated as absent.

Default implementations are conservative: unsupported operations reject
the request and no secure channel is considered open. This lets simple
bring-up Security Domains implement only the methods that are currently
needed while keeping the authority boundary explicit.

Identity and deployment role are not decided by the Rustlet itself.
In particular, the AID under which a Security Domain instance is known,
its \texttt{parent\_sd\_aid}, and whether that instance is the bootstrap root
Security Domain are kernel-side deployment properties rather than Rustlet
callbacks.

The binary relay used by the kernel is the generated
\texttt{sddispatch} entry in the Security Domain vtable. It is not a
second APDU transport. It is a privileged framework call from the kernel
to the selected Security Domain instance. The ABI carried through
\texttt{sddispatch} is intentionally separate from ordinary APDU
dispatch so that management decisions, secure-channel state, secure
messaging transforms, and
application APDU processing remain distinct.

The present implementation already validates that a Rustlet Security
Domain can be installed, selected, and reached through the generated
Security Domain vtable. The current \texttt{sddispatch} framework call
already routes the three main families of dedicated Security Domain
operations:
\begin{itemize}
\item \textbf{management}: load and installation authorization; deletion and
  data access; applet management; selectable-state decisions;
\item \textbf{secure channel}: profile support, session initialization and
  authentication, security-level inspection, open-state and MAC-length
  queries, and session reset;
\item \textbf{secure messaging}: \texttt{unwrap\_command()} and
  \texttt{wrap\_response()}.
\end{itemize}

The present reference Rustlet Security Domain already demonstrates this
model with a user-land SCP03 authority supporting \texttt{S8} and
\texttt{S16}, as well as protected exchanges using
\texttt{C-MAC/R-MAC} and \texttt{C-ENC/R-ENC}. The architectural point
is therefore the following: application behavior, management authority,
secure-channel establishment, secure messaging, and key-related policy
belong to distinct functional blocks, but they are exposed to the
developer through one coherent Rustlet-side Security Domain framework.

\section{Security Model and Isolation Principle}

\subsection{Isolation Principle}

Abstractly, a Rustlet is machine code executed inside the secure element
under operating-system control. It is not trusted as kernel code and it
does not execute with unrestricted access to secure-element memory.

The current isolation principle is based on four distinct memory areas:
\begin{itemize}
\item the Gate Region;
\item the execution stack;
\item the data and heap area;
\item the executable code area.
\end{itemize}

The Gate Region contains the shared APDU exchange buffer together with
the small execution gates required by the current isolated-call model.
The execution stack is the stack used while the Rustlet runs. The data
and heap area contains writable Rustlet-owned memory. The executable
code area contains the Rustlet machine code mapped from flash.

Physically, the execution stack and writable data belong to one
power-of-two RAM block allocated and aligned by the kernel allocator.
The stack occupies the lowest addresses and grows downward; relocated
data and the Rustlet heap immediately follow it at higher addresses.
The MPU exposes this complete block through one RW/XN region. A downward
stack overflow therefore leaves the permitted RAM region directly, without
requiring an internal guard gap or a second MPU region.

Any access outside these regions is denied by the MPU and results in a
fault. This is the essential isolation rule of the current model: a
Rustlet may access only the memory that has been explicitly assigned to
its execution context.

This isolation principle should be read together with the functional
model from Section~7.1. Because a Rustlet instance has a resident state,
and because installation and selection establish which instance is
executed, isolation is not only about code execution. It is also about
ensuring that one Rustlet instance cannot read or overwrite the state of
another instance, and cannot escape into kernel-owned memory.

Kernel services are reached through syscalls rather than through direct
calls into privileged kernel internals. In other words, a Rustlet does
not manipulate the operating system by linking against hidden kernel
functions. It requests services through explicit mediated entry points.

The effective realization of the MPU plan, the gate region, the return
path and the syscall trap handling is described in Chapter~8 (Reference
Implementation), in Section~\ref{sec:impl-rustlet-runtime}.

\subsection{Security Model}

The Rustlet security model should therefore be read at two levels. This
chapter defines the developer-facing contract: a Rustlet is isolated
application code, it owns only its assigned memory, it communicates with
the host through APDUs, and it reaches kernel services only through
mediated syscalls. The implementation details that make this contract
concrete are intentionally kept in Chapter~8, especially
Section~\ref{sec:impl-rustlet-runtime}: shared \texttt{RustletCtx} ABI,
selected-application descriptor, gate region, MPU programming,
syscall-return path, panic handling, recoverable faults, and serialized
state handoff.

This separation avoids mixing the application programming model with the
board- and kernel-specific mechanics. A Rustlet author should be able to
reason from the present chapter; a kernel or porting developer should use
Section~\ref{sec:impl-rustlet-runtime} to understand how the guarantee is
realized on a concrete target.

%--------------------------------------------------
